
Our results validate a four-component measurement framework for data quality in foundation model pretraining.

\subsubsection{Key Findings and Interpretation}

Deduplication as the strongest quality factor (r = 0.72) provides quantitative support for practitioner intuition. GPT-3, Gopher, and LLaMA all employ fuzzy deduplication, but prior to this work, no controlled experiments isolated its effect size relative to other quality dimensions. Our finding that dedup explains 52\% of information density variance---more than diversity (42\%), perplexity (34\%), or efficiency (28\%)---gives evidence-based priority ranking for curation resource allocation.

Information theory provides a mechanistic explanation: duplicate tokens carry zero marginal information. The second occurrence of identical text cannot teach a language model anything new, yet it consumes full compute cost during training. A corpus with 40\% duplicates wastes 40\% of gradient updates on zero-information examples. Removing them increases the effective information density per token.

Domain diversity as second-strongest (r = 0.65) validates The Pile's emphasis on multi-domain curation, but our quantitative comparison reveals an interesting hierarchy: deduplication > diversity. This suggests that within-domain redundancy removal matters more than cross-domain breadth.

Modest composite advantage (+8\% over best single metric) indicates quality dimensions are approximately linear. This has two implications: (1) production systems can use simple weighted averages (no need for neural quality scorers), and (2) strong non-linear interactions between quality factors are absent.

Reproducibility (CV < 10\%) addresses a critical practical concern: is Q(D) measurement stable enough for production use? Data quality metrics that vary wildly between samples would be useless for scaling law integration. Our low variance (average CV = 5.3\%) confirms Q(D) captures robust, replicable signals.

\subsubsection{Limitations}

\textbf{Domain Scope: Web Text Only}

We validated Q(D) on C4 (web-crawled English text), the most common pretraining corpus. Generalization to other domains---code, scientific text, multilingual---is untested. This is a principled boundary condition.

Why might Q(D) transfer or fail to transfer?

\begin{itemize}
\item Likely to transfer: Deduplication (redundancy is universal), token efficiency (boilerplate exists in all domains).
\item May require recalibration: Domain diversity (8 web categories vs 50 programming languages), perplexity (code has different syntactic patterns than prose).
\end{itemize}

Future work should test Q(D) on code and scientific corpora, expecting component weights to differ but correlation structure to hold.

\textbf{Scale Limitation: Correlation, Not Causation}

Our results validate correlation between Q(D) and information density. They do not prove that curation increases density (causality) or that higher density reduces compute requirements (efficiency claim). Without causal validation, our contribution is a measurement tool, not a complete scaling law.

\textbf{Model Scale: GPT-2 Small Only}

We used GPT-2 small (125M params) for perplexity scoring. Scale-invariance---whether Q(D) component weights hold at 1B-100B parameters---is untested. Two possibilities:

\begin{enumerate}
\item Universal Q(D): Same weights work across 125M-100B (ideal case)
\item Scale-dependent recalibration: Larger models weight diversity higher
\end{enumerate}

Literature provides weak evidence for (2): GPT-3 (175B) used more aggressive domain mixing than GPT-2 (1.5B). But this is correlation, not controlled experiment.

\subsubsection{Broader Impact}

\textbf{Positive impacts:}

\begin{enumerate}
\item Cost reduction: Q(D) measurement enables data-driven curation investment decisions.
\item Environmental benefit: Reducing redundancy means less wasted compute, lower carbon footprint.
\item Democratization: Smaller labs can achieve better performance with optimized data.
\end{enumerate}

\textbf{Negative impacts:}

\begin{enumerate}
\item Bias amplification risk: Q(D) optimizes for ''information density'' without fairness constraints. If deduplication removes minority-group text at higher rates, maximizing Q(D) could amplify demographic biases.
\item Measurement-target mismatch: Optimizing perplexity score (GPT-2-based) may overfit to GPT-2's specific biases.
\end{enumerate}

\textbf{Mitigation:} Future work should integrate fairness-aware Q(D) extensions---e.g., stratified deduplication (equal dedup rates across demographic groups), or multi-model perplexity consensus.

\subsubsection{Comparison to Prior Work}

Chinchilla provides L(N, D, C) with compute-optimal N-D ratios. Our Q(D) extends this toward L(N, D, Q(D), C). Chinchilla treats quality as constant; we measure it.

GPT-3 uses fuzzy dedup heuristically. We validate it quantitatively (r=0.72, strongest component).

The Pile emphasizes diversity (22 sources). We show dedup > diversity (r=0.72 vs 0.65) for web text---diversity matters, but redundancy removal matters more.

Data pruning methods remove low-value examples via model-based scoring. Our Q(D) is model-agnostic (computable before training), enabling pretraining-from-scratch use cases.

Our unique contribution is controlled experimental validation of quality metrics via information density, bridging qualitative practitioner knowledge and quantitative scaling law formalism.

\subsubsection{Future Directions}

Future work should:
\begin{itemize}
\item Test Q(D) on code and scientific corpora
\item Validate whether curation causally increases density (requires training experiments at scale)
\item Map compute-quality tradeoff curves (test ''does 20\% Q(D) improvement reduce compute?'')
\item Test scale-invariance across 1B-100B parameters
\end{itemize}

Long-term vision:
\begin{itemize}
\item Unified scaling law: L(N, D, Q(D), C) with fitted power-law coefficients
\item Pareto-optimal resource allocation: Solver for ''given budget, what (N, D, Q(D), curation\_cost) maximizes performance?''
\end{itemize}

\subsubsection{Positioning as Measurement Contribution}

We emphasize: this is a tool paper, not a solved scaling law. We provide a validated Q(D) measurement framework (r=0.78, CV<10\%) and evidence-based curation priority ranking. We do not provide causal proof that curation increases density or compute efficiency quantification. The contribution is making data quality measurable---the necessary first step before optimizing it.
